% =====================================================================
%  Lezione 5 — Giochi coalizionali e concetti di soluzione
%  Materiale: Lesson_5_Coalitional_Game_Theory_I.pdf,
%             Lab_5-Coalitional_Game_Theory_Core_Shapley.pdf,
%             CoalitionalGames_Frozensets_blank.ipynb
% =====================================================================
\chapter{Giochi coalizionali}

\parteteoria

Il percorso della lezione:
\begin{enumerate}[label=\arabic*.]
  \item introduzione ai giochi coalizionali;
  \item come rappresentarli;
  \item quali coalizioni si formano;
  \item come dividere il payoff: \kw{stabilità} (core) e \kw{equità} (nucleolo, valore di
        Shapley).
\end{enumerate}

% ---------------------------------------------------------------------
\section{Dai giochi non cooperativi a quelli coalizionali}

\begin{itemize}
  \item I giochi non cooperativi modellano giocatori egoisti e in competizione.
  \item Giochi come il dilemma del prigioniero o l'attraversamento mostrano che la
        competizione può portare a esiti subottimi.
  \item Gli equilibri correlati modellano (e rendono possibile) la cooperazione, ma
        introducono un dispositivo di coordinamento, un \emph{mediatore}.
  \item Nella realtà, però, le persone formano spontaneamente \kw{coalizioni} e prendono
        \kw{accordi vincolanti}.
\end{itemize}

La situazione tipica che si modella con un gioco coalizionale è questa:
\begin{itemize}
  \item i giocatori possono trarre vantaggio dal formare coalizioni;
  \item a ogni coalizione è associato un \kw{valore};
  \item il valore totale va \kw{distribuito} tra i giocatori.
\end{itemize}

Servono due scelte di modellazione:
\begin{center}
\renewcommand{\arraystretch}{1.3}
\begin{tabularx}{\linewidth}{@{}l X X@{}}
\toprule
 & Domanda & Modello \\
\midrule
Coalizioni & il risultato di una coalizione è indipendente dalle altre? &
 sì $\to$ coalizioni singole; no $\to$ più coalizioni insieme (giochi in funzione di
 partizione) \\
Valore & l'utilità è trasferibile? & sì $\to$ un solo numero per coalizione (TU);
 no $\to$ un insieme di vettori di utilità (NTU) \\
\bottomrule
\end{tabularx}
\end{center}

\begin{definizione}[Utilità trasferibile e non trasferibile]
L'utilità è \kw{trasferibile} (TU) se è la stessa per tutti i membri della coalizione e si
può ridistribuire tra loro: i payoff sono monetari (o convertibili in denaro), i giocatori
possono farsi pagamenti laterali, e la ridistribuzione non cambia il valore totale.
Esempi: dividere profitti, ripartire costi, assegnare seggi in parlamento.

L'utilità è \kw{non trasferibile} (NTU) se ogni membro ha la propria e la coalizione non
può convertirla o ridistribuirla liberamente. Esempi: il prestigio di un articolo scientifico
scritto insieme, i benefici personali di una collaborazione, i vantaggi legati a indicatori
di prestazione individuali.
\end{definizione}

\begin{attenzione}[denaro non vuol dire utilità trasferibile]
Il notebook lo dice bene: poter trasferire \emph{denaro} non implica che sia trasferibile
l'\emph{utilità}. Un giocatore ricco e uno povero possono trarre utilità molto diverse dalla
stessa somma (lezione 2: la funzione di Bernoulli è concava). Assumere TU significa assumere
che tutti valutino il denaro allo stesso modo, cioè utilità lineari nel denaro.
\end{attenzione}

\section{Giochi in funzione caratteristica}

\begin{definizione}[Gioco in funzione caratteristica]
Un gioco in funzione caratteristica è una coppia $(N, v)$ dove:
\begin{itemize}
  \item $N = \{1,\dots,n\}$ è l'insieme dei giocatori, detto anche \kw{grande coalizione};
  \item $v : 2^N \to \R$ è la \kw{funzione caratteristica}: per ogni sottoinsieme $C$ di
        giocatori, $v(C)$ è quanto possono guadagnare i membri di $C$ lavorando insieme.
        L'utilità $v(C)$ è trasferibile.
\end{itemize}
Di solito si assume che $v$ sia
\begin{itemize}
  \item \kw{normalizzata}: $v(\emptyset) = 0$;
  \item \kw{non negativa}: $v(C) \ge 0$ per ogni $C \subseteq N$;
  \item \kw{monotona}: $v(C) \le v(D)$ se $C \subseteq D$.
\end{itemize}
\end{definizione}

\begin{esempio}[il gelato]
Tre persone vogliono comprare del gelato e possono mettere insieme i soldi.
\begin{center}
\renewcommand{\arraystretch}{1.2}
\begin{tabular}{lc@{\qquad\qquad}lc}
\toprule
Giocatore & Soldi & Vaschetta & Prezzo \\
\midrule
A & \$6 & 500 g & \$7 \\
B & \$3 & 750 g & \$9 \\
C & \$3 & 1000 g & \$11 \\
\bottomrule
\end{tabular}
\end{center}
Il valore di una coalizione è la quantità massima di gelato che può comprare:
\[
\begin{aligned}
  v(\emptyset) &= v(\{A\}) = v(\{B\}) = v(\{C\}) = 0
    &&\text{(nessuno ha \$7 da solo)}\\
  v(\{A,B\}) &= v(\{A,C\}) = 750
    &&\text{(\$9: la vaschetta da 750 g)}\\
  v(\{B,C\}) &= 0 &&\text{(\$6: non basta per nulla)}\\
  v(\{A,B,C\}) &= 1000 &&\text{(\$12: la vaschetta da 1000 g).}
\end{aligned}
\]
\end{esempio}

\paragraph{Oltre la funzione caratteristica.}
I giochi in funzione caratteristica modellano utilità trasferibile e risultati di una
coalizione indipendenti da ciò che fanno gli altri. Se i risultati dipendono dalle altre
coalizioni si usano i \kw{giochi in funzione di partizione}, con valori $v(C, P)$ dove $C$
è una coalizione della partizione $P$; se l'utilità non è trasferibile si usano i
\kw{giochi NTU}, dove a ogni coalizione è associato un \emph{insieme} di vettori di utilità.
Ogni gioco TU si può rappresentare come un gioco NTU con un continuo di azioni, quindi
\[
  \text{NTU} \supset \text{TU} \supset \text{giochi in funzione caratteristica};
\]
nella pratica si usano soprattutto questi ultimi, che spesso si chiamano semplicemente
``giochi TU''.

\paragraph{Le due domande.}
I giochi coalizionali servono a rispondere a due domande:
\begin{enumerate}
  \item quali coalizioni si formeranno?
  \item come si divide il payoff all'interno di una coalizione?
\end{enumerate}
Nel gioco del gelato la coalizione con il benessere più alto è $\{A,B,C\}$. Ma dividere in
parti uguali non è equo: $A$ è più decisivo di $B$ e $C$, perché senza di lui non si compra
nulla.

% ---------------------------------------------------------------------
\section{Struttura di coalizione}

\begin{definizione}[Generazione della struttura di coalizione]
Dato un gioco $G = (N,v)$ si cerca la \kw{struttura di coalizione} $CS(G)$, cioè una
\emph{partizione} di $N$ in coalizioni disgiunte, che massimizza il benessere sociale
$\sum_{C \in CS(G)} v(C)$.
\end{definizione}

\begin{attenzione}[partizione, non ricoprimento]
Le slide dicono che una struttura di coalizione è un ``ricoprimento'' (\en{cover}) di $N$.
Un ricoprimento può avere insiemi sovrapposti, e un giocatore che sta in due coalizioni
verrebbe contato due volte. Quello che si intende, e che le slide stesse scrivono più
avanti, è una \kw{partizione}: le coalizioni coprono $N$ e sono a due a due disgiunte.
\end{attenzione}

\begin{esempio}[quando non conviene la grande coalizione]
\[
\begin{aligned}
  v(\emptyset) &= 0, & v(\{A\}) &= 0, & v(\{B\}) &= 300, & v(\{C\}) &= 0,\\
  v(\{A,B\}) &= 750, & v(\{A,C\}) &= 750, & v(\{B,C\}) &= 400, & v(\{A,B,C\}) &= 1000.
\end{aligned}
\]
La struttura di coalizione che massimizza il benessere è $\{\{A,C\}, \{B\}\}$, con
$750 + 300 = 1050 > 1000$: qui la grande coalizione \emph{non} è la scelta migliore.
\end{esempio}

\subsection{Il problema è difficile}

Le partizioni di un insieme di $n$ elementi sono il \kw{numero di Bell} $B_n$, che soddisfa
\[
  \left(\frac{n}{4}\right)^{n/2} \le B_n < n^n .
\]
Enumerarle tutte è impraticabile. Gli algoritmi usati sono la programmazione dinamica, il
grafo delle partizioni e le euristiche per giochi strutturati.

\subsection{Programmazione dinamica}

Sia $v^*(S)$ il massimo benessere ottenibile con una struttura di coalizione dei giocatori
di $S$, e $CS^*(S)$ la struttura che lo realizza: $v(CS^*(S)) = v^*(S)$.

\begin{teorema}[Equazione della programmazione dinamica]
Per ogni $C \subseteq N$:
\[
  v^*(C) = \max\Bigl\{ v(C),\ \max_{\substack{C' \cup C'' = C\\ C' \cap C'' = \emptyset}}
            \bigl(v^*(C') + v^*(C'')\bigr) \Bigr\}.
\]
\end{teorema}
O la coalizione $C$ resta intera, o la si divide in due parti, ciascuna delle quali va poi
organizzata nel modo migliore.

\paragraph{Dimostrazione.}
($\ge$) $v^*(C) \ge v(C)$, perché $\{C\}$ è una struttura di coalizione di $C$; e
$v^*(C) \ge v^*(C') + v^*(C'')$, perché unendo le strutture ottime delle due parti si ottiene
una struttura di $C$.
($\le$) Se $CS^*(C) = \{C\}$ allora $v^*(C) = v(C)$. Altrimenti sia $C'$ una coalizione di
$CS^*(C)$ e $C'' = C \setminus C'$. Il resto, $CS^*(C) \setminus \{C'\}$, è una struttura
di coalizione di $C''$, quindi $v^*(C'') \ge v^*(C) - v(C')$; inoltre
$v^*(C') \ge v(C')$. Sommando, $v^*(C') + v^*(C'') \ge v^*(C)$. $\square$

\begin{esempio}[la tabella della programmazione dinamica]
Con la funzione caratteristica dell'esempio precedente, procedendo per dimensione crescente:
\begin{center}
\renewcommand{\arraystretch}{1.25}
\begin{tabular}{lll}
\toprule
$C$ & $v^*(C)$ & $CS^*(C)$ \\
\midrule
$\{A\}$ & $0$ & $\{A\}$ \\
$\{B\}$ & $300$ & $\{B\}$ \\
$\{C\}$ & $0$ & $\{C\}$ \\
$\{A,B\}$ & $\max(v^*(A)+v^*(B),\ v(AB)) = \max(300, 750) = 750$ & $\{A,B\}$ \\
$\{A,C\}$ & $\max(0, 750) = 750$ & $\{A,C\}$ \\
$\{B,C\}$ & $\max(300, 400) = 400$ & $\{B,C\}$ \\
$\{A,B,C\}$ & $\max\bigl(1000,\ 750 + 0,\ 750 + 300,\ 400 + 0\bigr) = 1050$ & $\{\{A,C\},\{B\}\}$ \\
\bottomrule
\end{tabular}
\end{center}
Per $\{A,B,C\}$ le divisioni in due parti sono tre: $\{A,B\}|\{C\}$, $\{A,C\}|\{B\}$,
$\{B,C\}|\{A\}$. La migliore è la seconda.
\end{esempio}

\begin{attenzione}[un errore nella tabella delle slide]
Le slide scrivono $v^*(\{B,C\}) = \max(300, 400) = 750$ e riportano 750 in tabella. Il
massimo è 400. L'errore non cambia il risultato finale (1050 viene dalla divisione
$\{A,C\}|\{B\}$, che non usa $v^*(\{B,C\})$), ma nel passo successivo le slide elencano due
volte la divisione $\{A,B\}|\{C\}$ e mai $\{B,C\}|\{A\}$, proprio quella che userebbe il
valore sbagliato. Il codice del laboratorio conferma la tabella corretta.
\end{attenzione}

L'approccio è \kw{esatto}, ma il tempo è $O(3^n)$: non è praticabile per $n \ge 20$.

\begin{extra}[da dove viene $3^n$]
Si considerano tutte le coppie (coalizione $C$, sua parte $C'$). Ogni giocatore, in una
coppia, può stare in $C'$, in $C \setminus C'$ o fuori da $C$: tre possibilità per ciascuno
degli $n$ giocatori, cioè $3^n$ coppie. Con $n = 20$ sono circa $3{,}5$ miliardi.
\end{extra}

\paragraph{Grafo delle partizioni.}
Le partizioni si possono disporre su livelli, dal livello 1 (la sola grande coalizione) al
livello $n$ (tutti singoli); due partizioni sono collegate se una si ottiene dall'altra
fondendo due coalizioni. Le slide mostrano il grafo per 4 giocatori senza svilupparlo; l'idea
degli algoritmi che lo usano è esplorare per primi pochi livelli scelti bene e fermarsi
quando si ha una garanzia sufficiente sulla distanza dall'ottimo.

\paragraph{In pratica.}
Si usano euristiche e rappresentazioni compatte della funzione caratteristica: funzioni
caratteristiche analitiche, algoritmi di clustering (formazione delle coalizioni per
similarità o distanza, partizionamento di grafi), metaeuristiche (algoritmi genetici,
\en{simulated annealing}, \ldots).

% ---------------------------------------------------------------------
\section{Dividere il payoff}

\begin{definizione}[Esito di un gioco TU]
Un \kw{esito} di un gioco TU $G = (N,v)$ è una coppia $(CS, x)$ dove:
\begin{itemize}
  \item $CS = (C_1, \dots, C_k)$ è una struttura di coalizione, cioè una partizione di $N$:
        $\bigcup_i C_i = N$, $C_i \cap C_j = \emptyset$ per $i \ne j$;
  \item $x = (x_1,\dots,x_n)$ è un \kw{vettore di payoff} che distribuisce il valore di ogni
        coalizione di $CS$: $\sum_{i \in C} x_i = v(C)$ per ogni $C \in CS$
        (\kw{efficienza}).
\end{itemize}
Un esito è un'\kw{imputazione} se soddisfa la \kw{razionalità individuale}:
$x_i \ge v(\{i\})$ per ogni $i \in N$.
\end{definizione}

\begin{esempio}[esiti]
Con $v(\{1,2,3\}) = 9$ e $v(\{4,5\}) = 4$:
\begin{itemize}
  \item $\bigl((\{1,2,3\},\{4,5\}),\ (3,3,3,3,1)\bigr)$ è un esito: $3+3+3 = 9$, $3+1 = 4$;
  \item $\bigl((\{1,2,3\},\{4,5\}),\ (2,3,2,3,3)\bigr)$ \emph{non} lo è: la prima
        coalizione distribuisce $7$ e la seconda $6$. Non sono ammessi trasferimenti
        \emph{tra} coalizioni.
\end{itemize}
\end{esempio}

\subsection{Superadditività}

\begin{definizione}[Gioco superadditivo]
Un gioco $(N,v)$ è \kw{superadditivo} se
\[
  v(C_1 \cup C_2) \ge v(C_1) + v(C_2)
  \qquad \forall C_1, C_2 \subset N \text{ con } C_1 \cap C_2 = \emptyset .
\]
\end{definizione}

Il gioco del gelato è superadditivo: per esempio $v(\{A,B\}) = 750 \ge v(\{A\}) + v(\{B\}) = 0$,
$v(\{B,C\}) = 0 \ge 0 + 0$, $v(\{A,B,C\}) = 1000 \ge 750 + 0$.

Nei giochi superadditivi due coalizioni possono sempre fondersi senza perdere nulla: si può
assumere che si formi la \kw{grande coalizione}. Non c'è il problema della struttura di
coalizione e la notazione si semplifica: un esito è solo un vettore
$x = (x_1,\dots,x_n)$ con $\sum_{i\in N} x_i = v(N)$.

\begin{osservazione}[un'avvertenza delle slide]
Molti articoli definiscono gli esiti in questo modo, come semplici vettori con
$\sum x_i = v(N)$, \emph{anche} quando il gioco non è superadditivo. Va tenuto presente
quando si leggono: nel gioco della programmazione dinamica, per esempio, dividere $v(N)=1000$
sarebbe sbagliato, perché la struttura migliore produce $1050$.
\end{osservazione}

Da qui in poi si considerano solo giochi TU superadditivi.

% ---------------------------------------------------------------------
\section{Stabilità: il core}

Nel gioco del gelato si deve dividere $1000$. Se la divisione è $(200, 200, 600)$, $A$ e $B$
possono comprare da soli la vaschetta da 750 g e dividersela a metà, $375$ a testa: meglio
di $200$. L'esito $(200,200,600)$ \kw{non è stabile}.

\begin{definizione}[Core]
Il \kw{core} di un gioco è l'insieme degli esiti stabili, quelli da cui nessuna coalizione
vuole sganciarsi:
\[
  \core(G) = \Bigl\{ x \ \Big|\ x_i \ge 0\ \forall i \in N,\
  \sum_{i \in N} x_i = v(N),\
  \sum_{i \in C} x_i \ge v(C)\ \forall C \subseteq N \Bigr\}.
\]
Ogni coalizione riceve almeno quanto guadagnerebbe da sola.
\end{definizione}

\begin{esempio}[esiti nel core]
\begin{itemize}
  \item Con $v(\{1,2,3\}) = 9$, $v(\{4,5\}) = 4$ e $v(\{2,4\}) = 7$, l'esito
        $(3,3,3,3,1)$ \emph{non} è nel core: i giocatori 2 e 4 ricevono $3 + 3 = 6 < 7$ e
        preferiscono formare la coalizione $\{2,4\}$.
  \item Nel gelato, $(200,200,600)$ non è nel core perché
        $v(\{A,B\}) = 750 > x_A + x_B = 400$.
  \item Nel gelato, $(500,250,250)$ \emph{è} nel core: $x_A + x_B = 750 \ge 750$,
        $x_A + x_C = 750 \ge 750$, $x_B + x_C = 500 \ge 0$. Nessuna deviazione dà di più.
\end{itemize}
\end{esempio}

\begin{osservazione}[il primo esempio]
L'esempio con $(3,3,3,3,1)$ usa una struttura di coalizione $\{\{1,2,3\},\{4,5\}\}$, mentre la
definizione di core appena data è quella per la grande coalizione. Il concetto è lo stesso
(nessuna coalizione deve poter fare meglio da sola), ma nella versione con strutture di
coalizione l'efficienza si chiede coalizione per coalizione.
\end{osservazione}

\begin{extra}[com'è fatto il core del gelato]
I vincoli che contano sono due. Da $x_A + x_B \ge 750$ e $x_A + x_B + x_C = 1000$ segue
$x_C \le 250$; analogamente $x_B \le 250$. Il core è quindi
\[
  \core = \{ x \ge 0 : x_B \le 250,\ x_C \le 250,\ x_A = 1000 - x_B - x_C \},
\]
e in particolare $x_A \ge 500$ in ogni esito stabile. Si vede bene il potere di $A$: se $B$
chiedesse più di 250, $A$ potrebbe accordarsi con $C$ e dividere 750 dando a entrambi più di
quanto ricevono; lo stesso vale a parti invertite. $B$ e $C$ sono in concorrenza tra loro
per essere il partner di $A$.
\end{extra}

\subsection{Due cattive notizie}

\begin{esempio}[il core può essere vuoto]
\[
  v(\{1\}) = v(\{2\}) = v(\{3\}) = 0,\quad
  v(\{1,2\}) = 7,\ v(\{1,3\}) = 6,\ v(\{2,3\}) = 5,\quad
  v(\{1,2,3\}) = 8.
\]
Per stare nel core serve stabilità, $x_1 + x_2 \ge 7$, $x_1 + x_3 \ge 6$, $x_2 + x_3 \ge 5$,
ed efficienza, $x_1 + x_2 + x_3 = 8$. Sommando le tre disuguaglianze:
$2(x_1+x_2+x_3) \ge 18$, cioè $x_1+x_2+x_3 \ge 9 > 8$. Impossibile: il core è vuoto. E il
gioco è superadditivo, quindi la superadditività non basta a garantire un core non vuoto.
\end{esempio}

\begin{esempio}[il core può contenere esiti iniqui]
Nel gelato $(1000, 0, 0)$ è nel core: $x_A + x_B = 1000 \ge 750$, $x_A + x_C = 1000 \ge 750$,
$x_B + x_C = 0 \ge 0$, e la somma è $1000$. $B$ e $C$ mettono metà dei soldi e non ricevono
nulla, ma nessuna coalizione può fare meglio da sola. \kw{Stabilità non vuol dire equità.}
\end{esempio}

% ---------------------------------------------------------------------
\section{Il nucleolo}

\begin{definizione}[Eccesso]
L'\kw{eccesso} di una coalizione $S$ rispetto all'esito $x$ è
\[
  e(S, x) = v(S) - \sum_{i \in S} x_i .
\]
Misura l'\kw{insoddisfazione} della coalizione: quanto potrebbe guadagnare da sola rispetto
a quello che riceve. Un esito è nel core se e solo se tutti gli eccessi sono $\le 0$.
\end{definizione}

\begin{esempio}[eccessi]
Con $v(\{i\}) = 0$, $v(\{i,j\}) = 1$ per ogni coppia e $v(\{1,2,3\}) = 3$:
\begin{itemize}
  \item $x = (0,0,3)$: $e(\{1,2\},x) = 1 - 0 = 1$, la coppia $\{1,2\}$ è insoddisfatta;
  \item $x = (1,2,0)$: $e(\{1,2\},x) = 1 - 3 = -2$, la coppia è più che soddisfatta.
\end{itemize}
\end{esempio}

\begin{definizione}[Nucleolo]
Sia $\theta(x)$ il \kw{vettore degli eccessi} di tutte le coalizioni (diverse da $\emptyset$
e da $N$), ordinato dal più grande al più piccolo. Il \kw{nucleolo} è l'imputazione
\[
  \eta \in X(G) \ :\ \nexists\, y \in X(G) \text{ con } \theta(y) \prec_{lex} \theta(\eta),
\]
dove $X(G)$ è l'insieme delle imputazioni e $\prec_{lex}$ l'ordine lessicografico.
\end{definizione}
In parole: è l'esito che rende \kw{lessicograficamente minimo} il vettore delle
insoddisfazioni. Prima si rende la coalizione più insoddisfatta il meno insoddisfatta
possibile; a parità, la seconda; e così via.

\begin{esempio}[nucleolo del gioco simmetrico]
Nel gioco precedente:
\begin{itemize}
  \item $x = (1,2,0)$: $\theta(x) = (0, 0, -1, -1, -2, -2)$;
  \item $\eta = (1,1,1)$: $\theta(\eta) = (-1,-1,-1,-1,-1,-1)$.
\end{itemize}
Il primo elemento di $\theta(\eta)$ è $-1 < 0$: $\eta$ è lessicograficamente migliore, ed è
il nucleolo (lo conferma il calcolo con l'algoritmo seguente).
\end{esempio}

\begin{extra}[il nucleolo esiste sempre ed è unico]
Il nucleolo esiste ed è unico per ogni gioco con almeno un'imputazione. E se il core non è
vuoto il nucleolo sta nel core: un esito del core ha tutti gli eccessi $\le 0$, quindi il
nucleolo, che minimizza il più grande, ne ha anch'esso il più grande $\le 0$. È per questo
che lo si considera la versione ``più centrale'' del core.
\end{extra}

\subsection{L'algoritmo di Kopelowitz}

Il vettore degli eccessi si minimizza lessicograficamente con una \emph{sequenza} di
programmi lineari. Il primo minimizza l'eccesso massimo $\epsilon_1$:
\[
LP_1:\quad
\begin{aligned}
  \min\ & \epsilon_1\\
  \text{s.t. }\ & v(S) - \textstyle\sum_{i \in S} x_i \le \epsilon_1
    && \forall S \subset N,\ S \notin W_0 = \{\emptyset\}\\
  & x_i \ge v(\{i\}),\quad \textstyle\sum_{i \in N} x_i = v(N).
\end{aligned}
\]
Le coalizioni che all'ottimo hanno eccesso esattamente $\epsilon_1^*$ formano l'insieme
$W_1$ e il loro eccesso viene \emph{fissato}. Il programma $k$-esimo minimizza il massimo
eccesso delle coalizioni rimaste:
\[
LP_k:\quad
\begin{aligned}
  \min\ & \epsilon_k\\
  \text{s.t. }\ & v(S) - \textstyle\sum_{i \in S} x_i = \epsilon_r^*
    && \forall S \in W_r,\ r = 1,\dots,k-1\\
  & v(S) - \textstyle\sum_{i \in S} x_i \le \epsilon_k
    && \forall S \notin W_0 \cup W_1 \cup \dots \cup W_{k-1}\\
  & x_i \ge v(\{i\}),\quad \textstyle\sum_{i \in N} x_i = v(N),
\end{aligned}
\]
con $W_r = \{S \mid e(S,x) = \epsilon_r^*\}$. Ci si ferma quando tutte le coalizioni hanno
l'eccesso fissato.

\begin{itemize}
  \item Ogni programma ha $O(2^n)$ vincoli, uno per coalizione.
  \item Nel caso peggiore si fissa una coalizione per volta: fino a $2^n$ iterazioni.
  \item In generale è \kw{molto poco pratico}.
\end{itemize}

\begin{attenzione}[quali coalizioni fissare]
``Le coalizioni con eccesso pari a $\epsilon_k^*$'' va inteso come ``le coalizioni il cui
eccesso vale $\epsilon_k^*$ in \emph{tutte} le soluzioni ottime del programma $k$''. Una
soluzione ottima qualsiasi può avere per caso altre coalizioni con eccesso
$\epsilon_k^*$ che in un'altra soluzione ottima potrebbero scendere: fissarle porterebbe a un
risultato sbagliato. Nell'implementazione del laboratorio lo si controlla con un programma
lineare in più per ciascuna coalizione candidata (si prova ad abbassarne l'eccesso restando
all'ottimo).
\end{attenzione}

\subsection{Il nucleolo può essere iniquo}

\begin{esempio}[un gioco in cui conta solo A]
\begin{gather*}
  v(\emptyset) = v(\{A\}) = v(\{B\}) = v(\{C\}) = 0,\\
  v(\{A,B\}) = v(\{A,C\}) = 750,\quad v(\{B,C\}) = 0,\quad v(\{A,B,C\}) = 750.
\end{gather*}
Il nucleolo è $\eta = (750, 0, 0)$: assegna tutto ad $A$, anche se almeno uno tra $B$ e $C$
è necessario per ottenere qualcosa. Anche qui stabilità non implica equità.
\end{esempio}

\begin{attenzione}[il vettore degli eccessi nelle slide]
Le slide riportano $\theta(\eta) = (0,0,0,0,0,0)$. Gli eccessi di $\eta = (750,0,0)$ sono:
$e(\{B\}) = e(\{C\}) = 0$, $e(\{A,B\}) = e(\{A,C\}) = 750 - 750 = 0$, $e(\{B,C\}) = 0$, ma
$e(\{A\}) = 0 - 750 = -750$. Quindi $\theta(\eta) = (0,0,0,0,0,-750)$. Il nucleolo resta
quello: il core di questo gioco contiene \emph{solo} $(750,0,0)$ (da $x_A + x_B \ge 750$ e
$x_A + x_B + x_C = 750$ segue $x_C = 0$, e lo stesso per $x_B$), e il nucleolo sta nel core.
\end{attenzione}

% ---------------------------------------------------------------------
\section{Equità: il valore di Shapley}

\textbf{Idea}: un esito equo dovrebbe ricompensare ogni giocatore in proporzione al
\kw{contributo} che porta agli altri. Se il giocatore $i$ entra nella coalizione $C$, il suo
\kw{contributo marginale} è il surplus che crea:
\[
  v(C \cup \{i\}) - v(C).
\]
Nel gioco precedente, per esempio, $v(\{A,B\}) - v(\{A\}) = 750$ e
$v(\{A,B\}) - v(\{B\}) = 750$: il contributo di un giocatore dipende da chi c'è già. Un
esito equo si ottiene premiando ciascuno con il suo contributo \emph{medio}.

\begin{definizione}[Valore di Shapley]
Il valore di Shapley considera tutti gli ordini in cui la grande coalizione si può formare e
dà a ciascuno la stessa importanza:
\[
  \phi(i, v) = \frac{1}{|N|!} \sum_{\pi \in \Pi_N}
     \Bigl( v\bigl(B(\pi, i) \cup \{i\}\bigr) - v\bigl(B(\pi, i)\bigr) \Bigr),
\]
dove $\Pi_N$ è l'insieme delle permutazioni di $N$ e $B(\pi,i)$ l'insieme dei giocatori che
precedono $i$ nella permutazione $\pi$.
\end{definizione}

\begin{esempio}[Shapley nel gioco del gelato]
Contributi marginali di $A$ nei sei ordini:
\begin{center}
\renewcommand{\arraystretch}{1.2}
\begin{tabular}{lll}
\toprule
Ordine & Contributo di $A$ & \\
\midrule
$A,B,C$ & $v(\{A\}) - v(\emptyset)$ & $= 0$ \\
$A,C,B$ & $v(\{A\}) - v(\emptyset)$ & $= 0$ \\
$B,A,C$ & $v(\{A,B\}) - v(\{B\})$ & $= 750$ \\
$B,C,A$ & $v(\{A,B,C\}) - v(\{B,C\})$ & $= 1000$ \\
$C,A,B$ & $v(\{A,C\}) - v(\{C\})$ & $= 750$ \\
$C,B,A$ & $v(\{A,B,C\}) - v(\{B,C\})$ & $= 1000$ \\
\bottomrule
\end{tabular}
\end{center}
\[
  \phi(A) = \tfrac16(0 + 0 + 750 + 1000 + 750 + 1000) = 583{,}\overline{3}.
\]
Per $B$ i contributi sono $750, 250, 0, 0, 250, 0$, quindi
$\phi(B) = \tfrac{1250}{6} = 208{,}\overline{3}$, e per simmetria
$\phi(C) = 208{,}\overline{3}$. Somma: $1000$.
\end{esempio}

\begin{esempio}[Shapley nel gioco in cui conta solo A]
\begin{center}
\renewcommand{\arraystretch}{1.2}
\begin{tabular}{lccc}
\toprule
Ordine & contributo di $A$ & contributo di $B$ & contributo di $C$ \\
\midrule
$A,B,C$ & 0   & 750 & 0   \\
$A,C,B$ & 0   & 0   & 750 \\
$B,A,C$ & 750 & 0   & 0   \\
$B,C,A$ & 750 & 0   & 0   \\
$C,A,B$ & 750 & 0   & 0   \\
$C,B,A$ & 750 & 0   & 0   \\
\midrule
Valore di Shapley & $500$ & $125$ & $125$ \\
\bottomrule
\end{tabular}
\end{center}
A differenza del nucleolo, $B$ e $C$ ricevono qualcosa: chi arriva subito dopo $A$ è quello
che ``sblocca'' il valore.
\end{esempio}

\begin{extra}[equo, ma non stabile]
$(500, 125, 125)$ non è nel core di questo gioco: il core contiene solo $(750,0,0)$, e infatti
$x_A + x_B = 625 < 750 = v(\{A,B\})$. $A$ e $B$ da soli prenderebbero di più. Equità e
stabilità sono richieste diverse e possono essere incompatibili. Nel gioco del gelato,
invece, il valore di Shapley $(583{,}3;\ 208{,}3;\ 208{,}3)$ \emph{è} nel core (rispetta
$x_B, x_C \le 250$). Il risultato è coerente con l'esempio delle tre imprese della lezione 1,
dove con valore 100 si otteneva $(66{,}67;\ 16{,}67;\ 16{,}67)$: è lo stesso gioco in scala.
\end{extra}

\subsection{Caratterizzazione assiomatica}

\begin{teorema}[Shapley]
Il valore di Shapley è l'\kw{unico} schema di ripartizione che soddisfa:
\begin{enumerate}
  \item \kw{efficienza}: $\phi_1 + \phi_2 + \dots + \phi_n = v(N)$;
  \item \kw{giocatore nullo}: se $v(C \cup \{i\}) = v(C)$ per ogni $C$, allora $\phi_i = 0$;
  \item \kw{simmetria}: se $i$ e $j$ sono intercambiabili
        ($v(C \cup \{i\}) = v(C \cup \{j\})$ per ogni $C$ che non li contiene), allora
        $\phi_i = \phi_j$;
  \item \kw{additività}: $\phi_i(G_1 + G_2) = \phi_i(G_1) + \phi_i(G_2)$.
\end{enumerate}
Inoltre si può calcolare con la formula sui sottoinsiemi:
\[
  \phi(i, v) = \sum_{\substack{C \subseteq N\\ i \in C}}
     \frac{(|N| - |C|)!\,(|C| - 1)!}{|N|!}\,\bigl( v(C) - v(C \setminus \{i\}) \bigr).
\]
\end{teorema}

\begin{osservazione}[da dove viene il coefficiente]
Negli ordini in cui $i$ entra trovando già esattamente la coalizione $C \setminus \{i\}$,
i $|C|-1$ predecessori possono stare in $(|C|-1)!$ ordini e i $|N|-|C|$ successori in
$(|N|-|C|)!$. Il coefficiente è la frazione di permutazioni in cui $i$ dà proprio il
contributo $v(C) - v(C\setminus\{i\})$. Si passa da $n!$ permutazioni a $2^{n-1}$
coalizioni contenenti $i$: è la formula da usare nel codice.
\end{osservazione}

\begin{attenzione}[``giocatore fittizio'' e ``giocatore nullo'']
Le slide chiamano l'assioma 2 \emph{dummy player} e chiedono $\phi_i = 0$. In molti testi
il \emph{dummy} è un giocatore che aggiunge sempre esattamente $v(\{i\})$, e riceve
$\phi_i = v(\{i\})$; quello che riceve zero è il \emph{null player}, che non aggiunge mai
nulla. Con $v(\{i\}) = 0$ le due nozioni coincidono. Le slide scrivono anche ``Shapely''
invece di Shapley.
\end{attenzione}

% =====================================================================
\partelab{Core e valore di Shapley}
% =====================================================================
\fonte{CoalitionalGames\_Frozensets\_blank.ipynb, Lab\_5-Coalitional\_Game\_Theory\_Core\_Shapley.pdf}

\textbf{Obiettivo}: rappresentare un gioco coalizionale in Python, verificare la
superadditività, controllare se un esito sta nel core e calcolare il valore di Shapley.
Il notebook contiene gli output attesi; tutti sono riprodotti. In più sono stati implementati
e verificati la programmazione dinamica per la struttura di coalizione e il nucleolo.

\subsection{Rappresentare la funzione caratteristica}

Si usa un dizionario indicizzato da \texttt{frozenset}: i \texttt{frozenset} sono insiemi
immutabili, quindi si possono usare come chiavi, e hanno tutte le operazioni degli insiemi
(appartenenza, inclusione, unione, intersezione).

\begin{lstlisting}
from itertools import combinations
from math import factorial

v = {
    frozenset(['A']): 0, frozenset(['B']): 0, frozenset(['C']): 0,
    frozenset(['A', 'B']): 750, frozenset(['A', 'C']): 750,
    frozenset(['B', 'C']): 0,
    frozenset(['A', 'B', 'C']): 1000,
}

def powerset(players):   # tutte le coalizioni NON vuote: 2^n - 1
    return [frozenset(c) for k in range(1, len(players) + 1)
            for c in combinations(players, k)]
\end{lstlisting}

\begin{attenzione}[la coalizione vuota non c'è]
Il dizionario del notebook non contiene $\emptyset$, e la \texttt{powerset} fornita lo
esclude (il commento dice ``$2^{n-1}$ frozen sets'', ma sono $2^n - 1$). Ogni formula che
usa $v(\emptyset)$, come la formula di Shapley con $C = \{i\}$, deve trattarlo a parte:
\texttt{v.get(S, 0)} invece di \texttt{v[S]}, altrimenti un \texttt{KeyError}.
\end{attenzione}

\subsection{Superadditività}

\begin{lstlisting}
def is_superadditive(v):
    for C1 in v:
        for C2 in v:
            if not (C1 & C2):                       # disgiunte
                if v.get(C1 | C2, 0) < v[C1] + v[C2]:
                    return False
    return True
\end{lstlisting}
Il gioco del gelato dà \texttt{True}; il gioco del notebook con $v(A) = 6$, $v(B) = 12$,
$v(A,B) = 12$ dà \texttt{False} (perché $12 < 6 + 12$). Entrambi come atteso.

\subsection{Core}

\begin{lstlisting}
def is_stable(x, v, tol=1e-9):
    return all(sum(x[i] for i in C) >= v[C] - tol for C in v)

def is_in_the_core(x, v, tol=1e-9):
    N = max(v)                                       # la grande coalizione
    return (all(xi >= -tol for xi in x.values())
            and abs(sum(x.values()) - v[N]) <= tol
            and is_stable(x, v, tol))
\end{lstlisting}

\begin{center}
\renewcommand{\arraystretch}{1.2}
\begin{tabular}{lcl}
\toprule
Esito & Nel core? & Perché \\
\midrule
$(333{,}3334;\ 333{,}3333;\ 333{,}3333)$ & no & $x_A + x_B \approx 666{,}7 < 750$ \\
$(500;\ 250;\ 250)$ & sì & tutti i vincoli rispettati \\
$(1000;\ 0;\ 0)$ & sì & stabile, anche se iniquo \\
\bottomrule
\end{tabular}
\end{center}

\begin{attenzione}[\texttt{max(v)} trova la grande coalizione, ma per un motivo sottile]
Il notebook ricava la grande coalizione con \texttt{max(characteristic\_function)}. Funziona,
ma non perché confronti le dimensioni: per i \texttt{frozenset} l'operatore \texttt{>}
significa ``sovrainsieme proprio''. \texttt{max} tiene il primo elemento e lo sostituisce
con ogni elemento che lo contiene strettamente; la grande coalizione contiene tutte le altre,
quindi appena la incontra non la lascia più. Se la grande coalizione mancasse dal dizionario,
\texttt{max} restituirebbe una coalizione qualsiasi senza segnalare nulla.
\end{attenzione}

\begin{osservazione}[perché non si calcola il core]
Il notebook si limita a \emph{verificare} se un esito sta nel core. Calcolarlo tutto
richiederebbe di trovare i vertici del poliedro definito dai vincoli, con una sequenza di
programmi lineari. Per sapere solo se il core è \emph{vuoto} basta però un programma
lineare: si cerca un $x$ qualsiasi che soddisfi i vincoli. Applicato al gioco con
$v(\{1,2\}) = 7$, $v(\{1,3\}) = 6$, $v(\{2,3\}) = 5$, $v(N) = 8$ il risolutore conferma che
non esiste.
\end{osservazione}

\subsection{Valore di Shapley}

\begin{lstlisting}
def shapley_value(player, v):
    N = max(v)
    n = len(N)
    phi = 0.0
    for C in v:
        if player in C:
            s = len(C)
            weight = factorial(n - s) * factorial(s - 1) / factorial(n)
            phi += weight * (v[C] - v.get(C - {player}, 0))
    return phi

def shapley(v):
    return {p: shapley_value(p, v) for p in max(v)}
\end{lstlisting}

Risultato sul gelato: $\{A: 583{,}33;\ B: 208{,}33;\ C: 208{,}33\}$, come atteso. Una seconda
implementazione con le permutazioni dà gli stessi valori.

\paragraph{Sottoinsiemi o permutazioni.}
Il notebook raccomanda la formula sui sottoinsiemi: $O(2^n)$ invece di $O(n!)$. Su giochi
casuali con $n$ giocatori, tempi misurati per calcolare il valore di \emph{tutti} i giocatori:
\begin{center}
\renewcommand{\arraystretch}{1.2}
\begin{tabular}{rrr}
\toprule
$n$ & Sottoinsiemi & Permutazioni \\
\midrule
6  & 0,1 ms & 3,8 ms \\
8  & 0,8 ms & 304 ms \\
9  & 1,9 ms & 3\,409 ms \\
10 & 5,5 ms & $3\,628\,800$ ordini, non eseguito \\
\bottomrule
\end{tabular}
\end{center}
I valori calcolati con i due metodi coincidono ogni volta. Il codice sopra scorre tutte le
coalizioni per ogni giocatore, quindi il costo totale è $n\cdot 2^n$: comunque esponenziale.
Per giochi con molti giocatori si usano approssimazioni (per esempio campionando permutazioni
a caso) o rappresentazioni compatte, come anticipa la fine del notebook.

\begin{attenzione}[``il valore di Shapley è individualmente razionale'']
Il notebook elenca tra le proprietà $\phi(i,v) \ge v(\{i\})$. Vale per i giochi
\emph{superadditivi}: lì ogni contributo marginale è
$v(C \cup \{i\}) - v(C) \ge v(\{i\})$, e quindi lo è anche la media. Senza superadditività
può fallire. Il notebook stesso ricorda subito dopo l'altro limite: il valore di Shapley
non è necessariamente nel core.
\end{attenzione}

\subsection{In più: struttura di coalizione e nucleolo}

\paragraph{Programmazione dinamica.}
Implementazione diretta dell'equazione della teoria, coalizione per coalizione in ordine di
dimensione:
\begin{lstlisting}
def coalition_structure_dp(v):
    N = max(v)
    best, struct = {}, {}
    for C in sorted(powerset(sorted(N)), key=len):
        best[C], struct[C] = v[C], [C]
        members = sorted(C)
        for k in range(1, len(C)):
            for part in combinations(members, k):
                C1 = frozenset(part); C2 = C - C1
                if min(C1) != min(C):          # ogni divisione una volta sola
                    continue
                if best[C1] + best[C2] > best[C]:
                    best[C] = best[C1] + best[C2]
                    struct[C] = struct[C1] + struct[C2]
    return best, struct
\end{lstlisting}
Sul gioco dell'esempio restituisce esattamente la tabella corretta: $v^*(\{B,C\}) = 400$ e
$v^*(N) = 1050$ con $CS^* = \{\{A,C\},\{B\}\}$.

\paragraph{Nucleolo.}
Implementato con l'algoritmo di Kopelowitz e \texttt{linprog}, fissando a ogni passo solo le
coalizioni il cui eccesso non si può abbassare. Risultati:
\begin{center}
\renewcommand{\arraystretch}{1.2}
\begin{tabular}{lll}
\toprule
Gioco & Nucleolo & Vettore degli eccessi \\
\midrule
Simmetrico ($v(ij) = 1$, $v(N) = 3$) & $(1, 1, 1)$ & $(-1,-1,-1,-1,-1,-1)$ \\
Conta solo $A$ ($v(N) = 750$) & $(750, 0, 0)$ & $(0,0,0,0,0,-750)$ \\
Gelato ($v(N) = 1000$) & $(750, 125, 125)$ & $(-125,-125,-125,-125,-250,-750)$ \\
\bottomrule
\end{tabular}
\end{center}
Nel gelato nucleolo e Shapley sono diversi: $(750; 125; 125)$ contro
$(583{,}3; 208{,}3; 208{,}3)$. Entrambi stanno nel core, ma il nucleolo ``bilancia'' le
insoddisfazioni (le coppie $\{A,B\}$ e $\{A,C\}$ e i singoli $B$ e $C$ hanno tutti eccesso
$-125$), mentre Shapley premia i contributi marginali.

% ---------------------------------------------------------------------
\begin{riepilogo}
\begin{itemize}
  \item Gioco in funzione caratteristica $(N,v)$: $v(C)$ è quanto guadagna la coalizione
        $C$; utilità trasferibile.
  \item Struttura di coalizione: partizione di $N$ che massimizza
        $\sum v(C)$. Le partizioni sono il numero di Bell; la programmazione dinamica la
        trova in $O(3^n)$.
  \item Superadditività: fondersi non fa mai perdere, quindi si forma la grande coalizione e
        un esito è un vettore $x$ con $\sum x_i = v(N)$.
  \item \textbf{Core}: esiti in cui ogni coalizione riceve almeno $v(C)$. Può essere vuoto e
        può contenere esiti iniqui.
  \item \textbf{Nucleolo}: minimizza lessicograficamente gli eccessi
        $e(S,x) = v(S) - x(S)$; esiste, è unico, sta nel core se questo non è vuoto; si
        calcola con una sequenza di programmi lineari (Kopelowitz). Può essere iniquo.
  \item \textbf{Shapley}: contributo marginale medio su tutti gli ordini di ingresso; unico
        a soddisfare efficienza, giocatore nullo, simmetria, additività. Equo ma non
        necessariamente stabile.
  \item \textbf{Lab}: \texttt{frozenset} come chiavi, superadditività, verifica del core,
        Shapley sui sottoinsiemi ($n\cdot2^n$ contro $n!$); in più programmazione dinamica e
        nucleolo.
\end{itemize}
\end{riepilogo}
